\documentclass[10pt,a4paper]{amsart}
\usepackage[margin=19mm]{geometry}
\usepackage{amsmath,amssymb,mathtools}
\usepackage[scr=rsfs,cal=euler]{mathalfa}
\usepackage{xfrac}
\usepackage[unicode,hidelinks,bookmarks=false]{hyperref}
\newcommand{\ie}{i.e.\ }
\newcommand{\cf}{cf.\ }
\newcommand{\resp}{respectively\ }
\newcommand{\Subsection}{\subsection}
\providecommand{\nobreakdash}{\nobreak\hbox{-}}
\setlength{\emergencystretch}{2em}
\multlinegap0pt
\usepackage{mathtools}
\usepackage{xcolor}
\usepackage[normalem]{ulem}
\usepackage{braket}
\usepackage{amssymb}
\usepackage{bm}
\usepackage{float}
\usepackage{dsfont}
\usepackage{algorithm}\usepackage{algpseudocode}
\usepackage[shortlabels]{enumitem}
\newtheorem{lemma}{Lemma}
\usepackage{cleveref}
\crefname{lemma}{Lemma}{Lemmas}
\newcommand{\RR}{\mathbb{R}}
\newcommand{\VV}{\mathcal{V}}
\newcommand{\bbH}{\mathbb{H}}
\DeclareMathOperator{\capacity}{cap}
\title{Convex optimization on moment polytopes: Hadamard~mirror~descent and efficient algorithms for quantum functionals and other tensor parameters}
\author{Mahmut Levent Do\u{g}an, Keiya Sakabe, Michael Walter}
\date{Source: arXiv:2609.06633v1 (CC BY 4.0)}
\begin{document}
\maketitle

\noindent\emph{Source and licence.} Convex optimization on moment polytopes: Hadamard~mirror~descent and efficient algorithms for quantum functionals and other tensor parameters. Version arXiv:2609.06633v1 (CC BY 4.0), distributed by arXiv under the Creative Commons Attribution 4.0 International licence, http://creativecommons.org/licenses/by/4.0/, as recorded in the arXiv licence metadata for this exact version and shown on the abstract page https://arxiv.org/abs/2609.06633v1. Full licence notice: \texttt{licenses/arxiv-2609.06633-CC-BY-4.0.txt}.\par\smallskip

\noindent\textit{Editorial scope.} Counted unit: the article's lemma lem:shiftingtrick (source0.tex lines 3997--4003). The article's complete original proof of it is delivered with it (source0.tex lines 4004--4031, one \texttt{proof} environment of 28 lines). No mathematics is written, repaired or re-derived: the statement and the proof are the article's own lines, in the article's own notation, and no retained line is substituted or re-flowed.  Retained uncounted as the source-local context the proof needs: lines 999--1001; lines 1849--1849; lines 2132--2139.  Nothing was omitted from any retained span, so the package carries no omission record.  No retained line contains a citation, so the wrapper carries no bibliography and no bibliography entry.  No retained line contains a figure environment, an image-inclusion command or drawing code, so no image is delivered and the package declares no asset dependency.  Editorial formatting only, in addition: an AMS \texttt{amsart} preamble that restates the article's own declarations and macros used by the retained lines, namely the environments \texttt{lemma} on separate counters, together with the standard packages those lines need.  The retained environments print in this document as Lemma 1 in order of appearance, whereas the article prints the same items with the numbering of its own full document.  The complete native source of the article as distributed in the arXiv e-print for this exact version is bundled unchanged as \texttt{sources/209556/source0.tex}.
\begin{lemma}\label{lem:asymptoticrays}
    Any geodesic ray $t \mapsto g^\dagger e^{tX}g$ (where $g \in G, X \in \bbH$) on $P$ is asymptotic to $t \mapsto e^{tu^\dagger Xu}$ for some~$u \in K$.
\end{lemma}

\subsection{Convex optimization on moment polytopes}\label{sub:convex_optimization_on_moment_polytopes}

\begin{align}\label{eq:log cap K equivariance}
    \log\capacity_X(u \cdot v)
% = \frac12 \inf_{x \in P} f_{u \cdot v}^X(x)
% = \frac12 \inf_{x \in P} f^{u^\dagger X u}_v(u^\dagger x u)
% = \frac12 \inf_{x \in P} f^{u^\dagger X u}_v(x)
= \log\capacity_{u^\dagger X u}(v)
\qquad (u \in K, X \in \bbH).
\end{align}

\begin{lemma}\label{lem:shiftingtrick}
In the setting of \cref{sub:convex_optimization_on_moment_polytopes}, let $v \in \VV\setminus\{0\}$ and $X \in \bbH$.
Then:
    \[
        \exists g \in G,\ \log\capacity_X(g\cdot v) > -\infty \iff \exists u \in K,\ \log\capacity_X(u \cdot v) > -\infty.
    \]
\end{lemma}

\begin{proof}
Only the forward direction needs proof.
Let $g \in G$.
By \cref{lem:asymptoticrays}, there exists $u \in K$ such that the geodesic ray $\gamma(t) \coloneqq g^\dagger e^{tX}g$ is asymptotic to~$t \mapsto e^{t u^\dagger Xu}$, i.e., both determine the same equivalence class.
Then we have
\begin{align*}
    b^X(x)
= b^\gamma(g^\dagger x g)
= b^{u^\dagger X u}(g^\dagger x g) + c_g \qquad (x \in P),
\end{align*}
where the former equality follows from the isometry $x \mapsto g^\dagger xg$, and the latter from the fact that the Busemann functions of equivalent geodesics differ by a constant, which we call~$c_g \in \RR$.
Hence,
\begin{align*}
    f_{g\cdot v}^X(x)
= f_{g\cdot v}(x) + b^X(x)
= f_v(g^\dagger x g) + b^{u^\dagger X u}(g^\dagger x g) + c_g
= f_v^{u^\dagger X u}(g^\dagger x g) + c_g
\qquad (x \in P).
\end{align*}
By dividing by two and taking the infimum over $x \in P$, we obtain
\begin{align*}
    \log\capacity_X(g \cdot v)
= \log\capacity_{u^\dagger X u}(v) + \tfrac{c_g}2
= \log\capacity_X(u \cdot v) + \tfrac{c_g}2,
\end{align*}
where the second equality follows from \cref{eq:log cap K equivariance}.
Thus, if the left-hand side is finite, then so is the right-hand side.
\end{proof}

\end{document}
